% EGA V, English translation by P. Blass, 1987 (cote 24): archive cote 24, site page 31 (PDF page 32).
% typescript. Draft compared: none.
% AI-transcribed by Claude Opus 5.5 (claude-opus-5-5, adaptive thinking, effort high), Message Batches API, batch ega24b (msgbatch_01RGJJeHdR96U16ZWFuUjtDx), 2026-11-13; not yet checked by eye.
\scanpage{24/31}
\mspage{31}

trouble yourself with this exercise at least mention it as a
remark.  To say also that we do not discriminate in 6.1 with
regard to hyperplane sections that have an excessive (extra)
dimension.
(From the planning point of view) it might be clearer to group
together all the connectedness questions (including 3.3 and 4.4)
in the same No. (or section).

\uline{Section 7}

\uline{Application to the construction of hyperplane sections and}
\uline{multisections of specified type.}
\uline{7.1}

Let us notice that if S = Spec(k) where  k  is an infinite
field then every non-empty open subset of $P^{v}$ contains a k-rational
point therefore in the notations of 4.1 if E (defined in terms of
a constructible property P) contains the generic point $\eta$ , it
contains a k-rational point therefore there exists a hyperplane
section of X (defined over k) having the property P. On the other
hand, S being again arbitrary it is immediate that for every $s\in S$
and for every point $\xi$ of the fiber $Pv_{s}$ rational over k(s) there
exists a section $\xi$ of $Pv$ on an open neighborhood U of s which
passes through $\xi_{0}$.  If now E is again defined as in 4.1 in terms of
a constructible property P and if we have (for example due to No.
5) the fact that E is open then if $\xi_{0} \in E$ then the section \S{} is a
section of E over U at least if we sufficiently shrink or diminish
U.  Therefore we may construct a hyperplane section $\underline{Y}_{\xi}$ of X over an
open neighborhood U of s such that for every $t \in U$ its fiber $Y_{\xi(t)}$
